\documentclass[10pt]{article}
\usepackage[letterpaper,margin=0.85in]{geometry}
\input{glyphtounicode}
\pdfgentounicode=1
\usepackage{amsmath,amssymb,amsthm}
\usepackage{array,booktabs,flafter}
\usepackage[colorlinks=true,linkcolor=blue,citecolor=blue,urlcolor=blue]{hyperref}
\hypersetup{pdftitle={New Hadamard orders above 10\textasciicircum 4},pdfauthor={Jeffery Kline}}
\newtheorem{theorem}{Theorem}
\renewcommand{\thetheorem}{\Alph{theorem}}
\newtheorem{lemma}{Lemma}
\newtheorem{proposition}[lemma]{Proposition}
\theoremstyle{remark}
\newtheorem*{remark}{Remark}
\emergencystretch=2em
\DeclareTextSymbolDefault{\textbullet}{OMS}
\DeclareTextSymbolDefault{\textsection}{OMS}

\newcommand{\tr}{^{\mathsf T}}
\newcommand{\QC}{\mathrm{QC}}

\title{New Hadamard orders above $10^4$}
\author{Jeffery Kline}
\date{November 24, 2026\\[3pt]\small Version 0.1.0 (release candidate)}

\begin{document}
\maketitle

\begin{abstract}
A Hadamard matrix of order $n$ is an $n\times n$ matrix with entries $\pm1$ and pairwise orthogonal rows. We construct Hadamard matrices of six orders: $159052=4\cdot17\cdot2339$, $824332=4\cdot206083$, $22425772=4\cdot31\cdot223\cdot811$, $25072=16\cdot1567$, $46192=16\cdot2887$ and $47968=32\cdot1499$. The first three come from a theorem: if a quaternary conference matrix of order $n+1$ exists with $n$ odd, and $q=2n^r-1$ is prime for some $r\ge1$, then a Hadamard matrix of order $4(2n^r+1)$ exists. Its proof combines Turyn's power construction for conference matrices, which we extend to non-Hermitian quaternary matrices, with Kumar's 2026 quaternary form of Yamada's construction of order $4(q+2)$; the inputs are quaternary conference matrices of orders $142$ and $322$ published by Zaitsev and by Kumar. No printed table covers these three orders, and none of the constructions listed in Section~\ref{sec:new} reaches them from previously known ingredients. The odd part $206083$ of the second is a prime $p\equiv3\pmod 4$, so product constructions do not apply to it. The last three come from a second construction, which places two dihedral Golay pairs of Kamali and Kharaghani, of lengths $m$ and $n$, on commuting tensor factors and gives a Hadamard matrix of order $16(m+n)$. Each of these three orders is half the smallest order recorded for its odd part in Kumar's 2026 table.
\end{abstract}

\section{Introduction}\label{sec:intro}
A Hadamard matrix of order $n$ is an $n\times n$ matrix $H$ with entries $\pm1$ and $HH\tr=nI$. Such matrices are conjectured to exist at every order divisible by~$4$. The 2026 survey of Egan, Kotsireas, London and \'O~Cath\'ain~\cite{EKLO} lists the orders up to $10^4$ at which none is known; every one is $4p$, $4kp$ with $k\le13$, or $8p$, for a prime $p\equiv3\pmod4$. For $p\equiv1\pmod4$, Miyamoto's theorem gives a Hadamard matrix of order $4p$ from one of order $p-1$~\cite{Mi91}. For $p\equiv3\pmod4$ no such recursion is known: the known matrices of order $4p$ come from Paley's constructions, which need $4p-1$ or $2p-1$ to be a prime power, from a computer search, or from one of a few families such as Yamada's construction of order $4(q+2)$ for a prime power $q\equiv1\pmod4$~\cite{Ya89}.

In 2026 Kumar~\cite{Ku26} gave a quaternary form of Yamada's construction: for a prime $q\equiv1\pmod 4$, any quaternary conference matrix of order $(q+3)/2$ (zero diagonal, off-diagonal entries in $\{\pm1,\pm i\}$, orthogonal rows) gives a Hadamard matrix of order $4(q+2)$. Yamada's versions need a real symmetric conference matrix or a skew Hadamard matrix of order $(q+3)/2$; the former cannot exist when $(q+1)/2$ is not a sum of two squares. Kumar found the required quaternary matrices by computer and settled the open orders $2572$, $3436$, $3820$, $4204$ and $6796$ this way.

\paragraph{What is new.} Turyn~\cite{Tu71} raised a real conference matrix of order $n+1$ to one of order $n^r+1$, and Glazyrin~\cite{Gl26} showed that the same construction works for Hermitian matrices with unimodular entries. We observe that it needs neither symmetry nor the Hermitian property (Lemma~\ref{lem:power}). Feeding the powers into Kumar's construction gives Theorem~\ref{thm:A}, and applying it to two quaternary conference matrices found by Zaitsev~\cite{Za26} and Kumar~\cite{Ku26} gives the three orders of Theorem~\ref{thm:B}. The power construction and the quaternary template are not new; the new contributions are the extension to non-Hermitian matrices, the combination, and the three orders.

The second construction (Section~\ref{sec:golay}) starts from Kamali and Kharaghani's dihedral Golay pairs~\cite{KK98}, whose entries lie in a signed group generated by $i$ and $s$ with $i^2=-1$, $s^2=1$ and $is=-si$. Kamali and Kharaghani combined one such pair with a binary Golay pair. Representing two pairs on independent tensor factors makes their coefficients commute, so that any two such pairs, or complex Golay pairs, can be combined (Theorem~\ref{thm:C}). Kharaghani's 1992 array~\cite{Kh92} has the same order formula, but its amicability hypothesis fails for these inputs. Theorem~\ref{thm:D} gives three orders that improve Kumar's table~\cite{Ku26}, which updates that of Cati and Pasechnik~\cite{CP24}; unlike the orders of Theorem~\ref{thm:B}, their odd parts lie within those tables.

\begin{theorem}\label{thm:A}
Let $n$ be odd and suppose there is a quaternary conference matrix of order $n+1$. If $r\ge1$ and $q=2n^r-1$ is prime, then there is a quaternary Hadamard matrix of order $2(q+2)$ and a Hadamard matrix of order $4(q+2)=4(2n^r+1)$.
\end{theorem}

\begin{theorem}\label{thm:B}
There are Hadamard matrices of orders
\[
\begin{aligned}159052&=4(2\cdot141^2+1)=4\cdot17\cdot2339,\\ 824332&=4(2\cdot321^2+1)=4\cdot206083,\\ 22425772&=4(2\cdot141^3+1)=4\cdot31\cdot223\cdot811.\end{aligned}
\]
\end{theorem}

\begin{remark}
None of these orders is reached by the constructions in Table~4 of~\cite{EKLO}, or by the 2026 families listed in Section~\ref{sec:new}, from the ingredient sets specified there.
\end{remark}

\begin{theorem}\label{thm:C}
If there are dihedral Golay pairs (for example, complex Golay pairs) of lengths $m$ and $n$, then there is a Hadamard matrix of order $16(m+n)$.
\end{theorem}

\begin{theorem}\label{thm:D}
There are Hadamard matrices of orders $25072=16\cdot1567$, $46192=16\cdot2887$ and $47968=32\cdot1499$. For each of these odd parts, the smallest order recorded in Kumar's table~\cite{Ku26}, which updates that of Cati and Pasechnik~\cite{CP24}, is twice as large.
\end{theorem}

For $r=1$, Theorem~\ref{thm:A} is Kumar's construction. Its value lies in the larger $r$: a quaternary conference matrix of order $n^r+1$ is then built from one of order $n+1$, so a single matrix found by search yields one Hadamard order for each $r$ with $2n^r-1$ prime. Section~\ref{sec:more} lists these exponents for the published matrices.

\section{Quaternary conference matrices and their powers}\label{sec:power}
A \emph{quaternary conference matrix} of order $c$, written $\QC(c)$, is a $c\times c$ matrix $Q$ with zero diagonal, entries in $\{1,-1,i,-i\}$ off the diagonal, and $QQ^*=(c-1)I$. Real symmetric and skew-symmetric conference matrices are examples, and so are the matrices $\bigl(\chi(y-x)\bigr)_{x,y\in\mathbb F_q}$, bordered, for a quartic character $\chi$ of a finite field of order $p'\equiv1\pmod4$. A \emph{quaternary Hadamard matrix} of order $c$ is a $c\times c$ matrix $K$ with entries in $\{\pm1,\pm i\}$ and $KK^*=cI$. Multiplying rows and columns by units brings $Q$ to the normal form
\[
Q=\begin{pmatrix}0&\mathbf 1\tr\\ \mathbf 1&Z\end{pmatrix},
\]
where the \emph{core} $Z$, of order $n=c-1$, has zero diagonal, unit entries elsewhere, zero row and column sums, and $ZZ^*=nI-J$. Conversely any such $Z$ is the core of a $\QC(n+1)$. Below, $I$ and $J$ are the identity and all-one matrices of the order required by the context (order $n$ unless indicated), and $\otimes$ is the Kronecker product.

\begin{lemma}[Turyn's power construction]\label{lem:power}
Let $Z$ be the core of a $\QC(n+1)$ and $r\ge2$. Then there is a core $Z_r$ of a $\QC(n^r+1)$. For $r=2$ one may take $Z_2=Z\otimes Z+I\otimes J-J\otimes I$. No symmetry of $Z$ is needed.
\end{lemma}

\begin{proof}
For $r=2$: the diagonal of $Z_2$ is $0+1-1=0$. At an off-diagonal position exactly one of the three terms is nonzero, and it is a unit. Since $ZJ=JZ=0$, the cross terms involving $Z\otimes Z$ vanish, the other two cross terms are each $-J\otimes J$, and
\[
Z_2Z_2^*=(nI-J)\otimes(nI-J)+nI\otimes J+nJ\otimes I-2J\otimes J=n^2I-J\otimes J .
\]
Row and column sums of $Z_2$ are zero because those of $Z$ are and $I\otimes J$, $J\otimes I$ have equal line sums. For odd $r\ge3$ we follow Turyn~\cite{Tu71} and Glazyrin~\cite[Theorem~5]{Gl26}. Call a word $S=S_1\cdots S_r$ in the letters $W,I,J$ a \emph{string} if, read cyclically, every $I$ is followed by $J$ and every $J$ is preceded by $I$, and put $G_S=S_1\otimes\cdots\otimes S_r$ with $W=Z$. Set $Z_r=\sum_S G_S$. Glazyrin's proof shows that the zero/unit pattern of $Z_r$ is right, and that two distinct strings have, at some position, $W$ in one and $J$ in the other. Here that gives $G_{S}G_{S'}^*=0$ for $S\ne S'$, because $ZJ=0$ and $JZ^*=(ZJ)^*=0$. Each $G_SG_S^*$ is the Kronecker product of factors $ZZ^*=nI-J$, $I$ and $J^2=nJ$, which are the factors of Glazyrin's computation with $W^2$ replaced by $WW^*$; his cancellation argument then gives $Z_rZ_r^*=n^rI-J$. Every string contains $W$ when $r$ is odd, so $G_SJ=JG_S=0$ and $Z_r$ has zero line sums. Finally, for $r=2^ar'$ with $r'$ odd, take the $r'$-th power and square it $a$ times.
\end{proof}

\begin{remark}
For real cores this is Turyn's theorem; Belevitch had the case $r=2$~\cite{GS67}. For Hermitian unimodular cores it is Lemma~4.1 and Theorem~5 of~\cite{Gl26}. The cores used in Section~\ref{sec:orders} are neither Hermitian nor skew-Hermitian. We checked Lemma~\ref{lem:power} exactly for the non-Hermitian quartic-character cores of orders $5$ ($r=3,5$) and $13$ ($r=3$).
\end{remark}

\begin{lemma}[Kumar~\cite{Ku26}, Theorem~3]\label{lem:template}
Let $q$ be a prime with $q\equiv1\pmod4$. If there is a $\QC((q+3)/2)$, then there is a quaternary Hadamard matrix $K$ of order $2(q+2)$, and
\[
H=\begin{pmatrix}K_{\mathrm{re}}+K_{\mathrm{im}}&K_{\mathrm{re}}-K_{\mathrm{im}}\\-(K_{\mathrm{re}}-K_{\mathrm{im}})&K_{\mathrm{re}}+K_{\mathrm{im}}\end{pmatrix},\qquad K=K_{\mathrm{re}}+iK_{\mathrm{im}},
\]
is a Hadamard matrix of order $4(q+2)$.
\end{lemma}

Kumar's matrix $K$ is a $2\times2$ array $\left(\begin{smallmatrix}X&Y\\Y^*&-X^*\end{smallmatrix}\right)$ whose blocks are bordered by a row and a column of units. The core $X'$ of $X$ is built from Yamada's symmetric circulants $M,N$ of order $(q+1)/2$ with $M^2+N^2=qI$~\cite{Ya89}, and the core of $Y$ is $Y'=J_2\otimes Z+P\otimes I$ with $P=\left(\begin{smallmatrix}1&-1\\-1&1\end{smallmatrix}\right)$ and $Z$ the core of the given $\QC((q+3)/2)$. The final step is Turyn's doubling~\cite{Tu70}. We do not repeat the proof; our code rebuilds $K$ and checked it exactly at orders $38$, $198$ and $1286$ (the last is Kumar's own instance at $q=641$).

\begin{proof}[Proof of Theorem~\ref{thm:A}]
For $r=1$ the conference matrix is given; for $r\ge2$ Lemma~\ref{lem:power} gives a $\QC(n^r+1)$. In both cases $n^r+1=(q+3)/2$. Since $n$ is odd, $q=2n^r-1\equiv1\pmod4$. Lemma~\ref{lem:template} applies.
\end{proof}

\section{The three orders}\label{sec:orders}
Theorem~\ref{thm:B} follows from Theorem~\ref{thm:A} applied to the two inputs below with the exponents of Table~\ref{tab:orders}; the comparison with known constructions is in Section~\ref{sec:new}.

\paragraph{Inputs.} Kumar's certificate \texttt{C1286.json}~\cite{Ku26} gives four sequences on $\mathbb Z_{161}$ which, by his zero-diagonal array (his Theorem~1), form a $\QC(322)$; we rebuilt it and checked $QQ^*=321I$ exactly. Zaitsev's quaternary Hadamard matrix of order $566$~\cite{Za26} contains the block $Y'$ described after Lemma~\ref{lem:template}, with $q=281$: its rows $2$--$283$ and columns $285$--$566$ form $Y'=J_2\otimes Z+P\otimes I$ exactly, so its upper left $141\times141$ corner minus $I$ is the core $Z$ of a $\QC(142)$. We checked both $KK^*=566I$ for Zaitsev's matrix and that $Z$ is a core, exactly. Neither core is Hermitian, skew-Hermitian or symmetric.

\paragraph{Construction.} Table~\ref{tab:orders} lists the three applications of Theorem~\ref{thm:A}. In each case $q$ is prime and $q\equiv1\pmod8$.

\begin{table}[h]
\centering
\begin{tabular}{rrrrrl}
\toprule
input & $r$ & $\QC$ order $n^r+1$ & $q=2n^r-1$ & order $4(q+2)$ & odd part\\
\midrule
$\QC(142)$ & 2 & 19882 & 39761 & 159052 & $39763=17\cdot2339$\\
$\QC(322)$ & 2 & 103042 & 206081 & 824332 & $206083$ (prime)\\
$\QC(142)$ & 3 & 2803222 & 5606441 & 22425772 & $5606443=31\cdot223\cdot811$\\
\bottomrule
\end{tabular}
\caption{The orders of Theorem~\ref{thm:B}.}\label{tab:orders}
\end{table}

\paragraph{Verification.} The matrices have orders up to $2.2\times10^7$, so we do not store them or form their Gram matrices. Instead the code checks every ingredient exactly and the final identity by random projection. (i) The input cores satisfy $ZZ^*=nI-J$ with zero line sums, exactly. (ii) The first rows $\alpha,\beta$ of Yamada's circulants $M,N$ at each $q$ have periodic autocorrelations summing to $q\delta$ (checked by a rounded FFT, and exactly at $43$ lags), are symmetric, and have the stated zero pattern. (iii) The power core satisfies $Z_rZ_r^*x=(n^rI-J)x$ for random vectors $x$ with entries in $\{\pm1,\pm i\}$. (iv) $KK^*x=K^*Kx=2(q+2)x$ for random vectors $x$ with entries in $\{\pm1,\pm i\}$ (three vectors each at $159052$ and $824332$, two at $22425772$), after checking that the coded $K^*$ is the adjoint of $K$, with $K$ and $K^*$ applied through their Kronecker and circulant structure. (v) Selected rows and columns of $K$ have unit entries. A wrong entry in $K$ or in an input makes a check of type (iii) or (iv) fail with probability at least $1/2$ per vector; deliberately corrupted inputs fail as expected. Every check except (iv) at $22425772$ runs in about a minute on one core; that one takes about four minutes and $2$\,GB of memory. Exactness of (ii)--(v) relies on rounding floating-point FFT output to Gaussian integers; the largest rounding residual observed was $2.4\times10^{-8}$.

\section{Comparison with known constructions}\label{sec:new}
No printed table of known Hadamard orders reaches this range: the survey~\cite{EKLO} stops at $10^4$, and the tables of Seberry and Yamada~\cite{SY20} and Cati and Pasechnik~\cite{CP24} cover odd parts up to $2999$. Seberry and Yamada describe a computer enumeration for odd parts below $40000$~\cite{SY92}, which would include $39763$, but their printed table stops far below it. We therefore tested each order against every construction in Table~4 of~\cite{EKLO} that can produce an order $4m$ with $m$ odd: Paley I and II, the Williamson array, the theorem of Cooper and Wallis (T-matrices and Williamson-type matrices are as in~\cite{SY20}), the Seberry--Yamada product with T-matrices (Paley and Miyamoto blocks), Miyamoto's constructions, Spence, Mathon, Seberry--Whiteman, and Yamada's two forms of $4(q+2)$; and against the 2026 families: Kumar's template itself with previously known cores, and the family $4(4q^2+3)$ of~\cite[Theorem~7]{EKLO}. Constructions such as Agaian's, Craigen--Seberry--Zhang, Craigen--Holzmann--Kharaghani and Kronecker products give orders divisible by~$8$ and do not apply.

The test over-approximates every ingredient: it treats T-matrices as known at every order up to $200$ and at all products of the known infinite families, Williamson-type matrices as known at every odd order below $2000$, Hadamard matrices as known at every order below $10^4$ outside the open list of~\cite{EKLO}, and conference matrices (symmetric, skew or quaternary) as known at every order up to $4000$, at every prime power plus one, and at the Mathon and Seberry--Whiteman orders. None of the three orders is reached even so. Turyn powers of real symmetric conference matrices are not in this list, and need not be: the three core orders are $141^2+1$, $321^2+1$ and $141^3+1$, neither $141$ nor $321$ is a sum of two squares, so no symmetric conference matrix of order $142$ or $322$ exists, and $141^3$ is not a sum of two squares either. At $159052$ the binding obstruction is that the Seberry--Yamada product with T-matrices of order $17$ would need a Hadamard matrix of order $4\cdot2339=9356$, which is itself open~\cite{EKLO}. At $824332$ the odd part is prime, so the product constructions would need T-matrices or Williamson-type matrices of order $206083$, which no known family supplies; Paley's constructions fail ($4\cdot206083-1=3\cdot274777$, $2\cdot206083-1=5\cdot13\cdot17\cdot373$), and Yamada's would need a real symmetric conference matrix of order $103042=321^2+1$; none is known, and squaring cannot produce one because $321$ is not a sum of two squares. The code reproduces the test (\texttt{code/openness.py}); positive controls include the orders $2572$, $5788$ and $9172$. We also searched GitHub, Zenodo, arXiv and the web for the three orders and their odd parts and found no construction (recorded in \texttt{audit/PRIOR-ART.md} of the release repository).

\section{Further orders}\label{sec:more}
Theorem~\ref{thm:A} applies to every known quaternary conference matrix of order $n+1$ with $n$ odd. Real symmetric conference matrices give nothing new: their Turyn powers are again real symmetric, so Yamada's theorem already applies. The interesting inputs are the non-real matrices found by search. Table~\ref{tab:more} lists, for the quaternary conference matrices underlying Zaitsev's and Kumar's quaternary Hadamard matrices, the exponents $2\le r\le16$ for which $2n^r-1$ is prime (a probable prime for large values). We have extracted and checked only the cores $n=141$ and $n=321$; the other rows assume that the corresponding matrices have the layout of Lemma~\ref{lem:template}, as Kumar states~\cite{Ku26}. We have not tested the resulting orders for newness, and $r=1$ returns the orders of the input matrices.

\begin{table}[h]
\centering
\begin{tabular}{rll}
\toprule
$n$ & source of $\QC(n+1)$ & $2\le r\le16$ with $2n^r-1$ prime\\
\midrule
117 & Zaitsev, order 470 & 5\\
141 & Zaitsev, order 566 & 2, 3, 14\\
201 & Zaitsev, order 806 & 4, 12\\
285 & Zaitsev, order 1142 & 3, 7, 10\\
321 & Kumar, order 1286 & 2, 6, 14\\
381 & Zaitsev, order 1526 & 5, 8, 9, 13\\
429 & Kumar, order 1718 & 4, 7, 10, 15\\
477 & Kumar, order 1910 & 10\\
525 & Kumar, order 2102 & 6, 7, 9\\
849 & Kumar, order 3398 & 3, 4\\
\bottomrule
\end{tabular}
\caption{Exponents $r\ge2$ for which Theorem~\ref{thm:A} applies, given the quaternary conference matrix of order $n+1$ (extracted and checked only for $n=141$ and $n=321$).}\label{tab:more}
\end{table}

Whether a given $n$ admits infinitely many such $r$ is a question about primes of the form $2n^r-1$, which is open. For $r\ge2$, Theorem~\ref{thm:A} does not reach any open order below $10^4$: that would need $2n^2+1<2500$, so $n\le35$; among odd $n$ that are not prime powers, $2n^2-1$ is prime only for $n=15$ and $n=21$, which give the known orders $1804$ and $3532$ (for prime powers $n$, Yamada's theorem already applies).

\section{Two dihedral Golay pairs}\label{sec:golay}
A \emph{dihedral Golay pair} of length $m$~\cite{KK98} is a pair $(a,b)$ of sequences of length $m$ with entries in $\{\pm1,\pm i,\pm s\}$, where $i^2=-1$, $s^2=1$ and $is=-si$, whose aperiodic autocorrelations are complementary and whose Laurent polynomials $a(z),b(z)$ commute with each other and with their adjoints. Complex Golay pairs are examples; for two complex Golay pairs Craigen, Holzmann and Kharaghani~\cite{CHK02} already give order $8(m+n)$, so Theorem~\ref{thm:C} adds something only when at least one pair is dihedral. Kamali and Kharaghani found such pairs of lengths $7$, $9$, $15$ and $19$ (we use the first three), and their Lemma~3 multiplies the length by any binary Golay length.

\begin{proof}[Proof of Theorem~\ref{thm:C}]
Represent $i$ by the rotation $\left(\begin{smallmatrix}0&-1\\1&0\end{smallmatrix}\right)$ and $s$ by $\mathrm{diag}(1,-1)$. Embed each coefficient $g$ of the first pair as $g\otimes I_2$ and each coefficient $h$ of the second pair as $I_2\otimes h$. The four resulting polynomials $a,b,c,d$, whose coefficients are $4\times4$ signed permutation matrices, and their adjoints commute pairwise: within a pair by hypothesis, and across pairs because the coefficient factors are independent. Let $N=m+n$, evaluate at the cyclic shift of order $N$, and set
\[
A=a+z^mc,\quad B=a-z^mc,\quad C=b+z^md,\quad D=b-z^md .
\]
The two aperiodic Golay identities survive the evaluation, so $AA\tr+BB\tr+CC\tr+DD\tr=4NI$, and $A,B,C,D$ and their transposes commute pairwise. Let $R$ be the reversal of order $N$ and $\mathcal J=R\otimes s\otimes s$. Then $\mathcal J\tr=\mathcal J$, $\mathcal J^2=I$, and $\mathcal JX\mathcal J=X\tr$ for $X=A,B,C,D$, because reversal inverts the shift and $sgs=g\tr$ for $g\in\{\pm1,\pm i,\pm s\}$. In Craigen's Goethals--Seidel array~\cite{Cr94}, as stated by Kamali and Kharaghani~\cite[Lemma~2]{KK98} and used in their Theorem~7, with $\mathcal J$ in place of their reflection,
\[
\begin{pmatrix}
A&-B\mathcal J&-C\mathcal J&-D\mathcal J\\
B\mathcal J&A&-D\tr\mathcal J&C\tr\mathcal J\\
C\mathcal J&D\tr\mathcal J&A&-B\tr\mathcal J\\
D\mathcal J&-C\tr\mathcal J&B\tr\mathcal J&A
\end{pmatrix},
\]
each diagonal block of the Gram matrix is $4NI$, and each off-diagonal block is a sum of commutators such as $[A,B]\mathcal J+[C,D]$, which vanish. Every entry of this $4N\times4N$ array is a signed $4\times4$ permutation matrix. Multiplying on the right by $I_{4N}\otimes H_4$, with $H_4$ a real Hadamard matrix of order $4$, gives a $\pm1$ matrix of order $16N$ with Gram matrix $16NI$.
\end{proof}

With the original two-dimensional representation of the coefficients the construction fails for the inputs below: the cross-commutators of the two pairs would have to vanish, and for the published seeds they do not. The tensor embedding is what makes the two pairs independent. Kharaghani's 1992 Theorem~9~\cite{Kh92} has the same order formula $4k(m+n)$ for block Golay inputs, but its Lemma~7 requires the coefficient blocks to be pairwise amicable, including two blocks on the same side, and the inputs used here violate that requirement.

\begin{proof}[Proof of Theorem~\ref{thm:D}]
Apply Theorem~\ref{thm:C} with the following lengths, each a published dihedral or complex Golay length multiplied by a binary Golay length by~\cite[Lemma~3]{KK98}: $7$ and $1560$ give $16\cdot1567=25072$; $7$ and $2880$ give $16\cdot2887=46192$; here $1560=15\cdot104$ and $2880=9\cdot320$ are dihedral lengths, and also complex Golay lengths $3\cdot520$ and $18\cdot160$, which the code uses; $7\cdot26=182$ and the complex Golay length $11\cdot256=2816$ give $16\cdot2998=47968$. Kumar's table~\cite{Ku26} records, for each odd part $m$, the least $t$ with a known Hadamard matrix of order $2^tm$; for $1567$, $2887$ and $1499$ these are $5$, $5$ and $6$, that is, the orders $50144$, $92384$ and $95936$.
\end{proof}

These three orders lower the least known power of $2$ for their odd parts; they do not settle the orders $4\cdot1567=6268$, $4\cdot2887=11548$ and $4\cdot1499=5996$, which remain open, nor the corresponding orders $8\cdot1567$, $8\cdot2887$ and $8\cdot1499$ (and $16\cdot1499$).

The three matrices have explicit entry formulas. The code checks the full Gram identity of each matrix exactly, by cyclic polynomial multiplication of the blocks, together with the input pairs, a dense control of order $160$ and a deliberately corrupted coefficient, which fails. A search of the sources available to us found earlier constructions of the comparison orders $50144$, $92384$ and $95936$, but none of the three new orders. Over all published inputs, Theorem~\ref{thm:C} improves no other entry of Kumar's table for odd parts up to $2999$. Above that range the theorem gives $94$ further orders up to $10^6$ outside the tables and the families we compared against. Earlier constructions exist for nine of them; for the other $85$, the smallest being $69232=16(7+4320)$, we located none. We do not claim these: no table covers them, and our comparison with earlier constructions there is incomplete.

\begin{remark}
Stopping the proof of Theorem~\ref{thm:C} before the final multiplication by $I_{4N}\otimes H_4$ leaves a $4N\times4N$ array $G$ of signed $4\times4$ permutation matrices with $GG\tr=4NI$. For any Hadamard matrix $F$ of order $4k$, $(G\otimes I_k)(I_{4N}\otimes F)$ is a Hadamard matrix of order $16k(m+n)$; this is the classical expansion of a signed-group matrix by a real Hadamard fill~\cite{Cr95}. With published inputs it gives $38$ orders up to $10^6$ outside the comparison above; earlier constructions exist for nine of them (recorded in the release repository), and priority is unresolved for the other $29$, the smallest being $125360$. We do not claim these either.
\end{remark}

\begin{remark}
Du and Jiang's Corollaries~3 and~5~\cite{DJ24} give a four-phase complementary quadruple of length $5(550+648)=5990$ and hence a Hadamard matrix of order $47920=16\cdot2995$. This lowers Kumar's recorded exponent for the odd part $2995$ from $5$ to $4$. It follows from their published results and is not a contribution of this note.
\end{remark}

\section{Role of AI}\label{sec:ai}
AI assisted the derivations, the computer verification, the literature search and the exposition under the author's direction. Process-separated AI reviews checked the proofs and certificates; they are not independent expert peer review, and no independent human expert has examined the proofs. The release repository records the sources, the reviews and their dispositions.

\begin{thebibliography}{99}
\bibitem{CP24} M.~Cati and D.~V.~Pasechnik, \href{https://arxiv.org/abs/2411.18897}{A database of constructions of Hadamard matrices}, arXiv:2411.18897v2, 2025.
\bibitem{Cr94} R.~Craigen, Complex Golay sequences, J.\ Combin.\ Math.\ Combin.\ Comput.\ \textbf{15} (1994), 161--169.
\bibitem{Cr95} R.~Craigen, Signed groups, sequences, and the asymptotic existence of Hadamard matrices, J.\ Combin.\ Theory Ser.\ A \textbf{71} (1995), 241--254.
\bibitem{CHK02} R.~Craigen, W.~Holzmann and H.~Kharaghani, Complex Golay sequences: structure and applications, Discrete Math.\ \textbf{252} (2002), 73--89.
\bibitem{DJ24} C.~Du and Y.~Jiang, \href{https://arxiv.org/abs/2401.15381}{Golay complementary sequences of arbitrary length and asymptotic existence of Hadamard matrices}, arXiv:2401.15381, 2024.
\bibitem{EKLO} R.~Egan, I.~S.~Kotsireas, S.~London and P.~\'O~Cath\'ain, \href{https://arxiv.org/abs/2610.27175}{Recent progress in the construction of Hadamard matrices}, arXiv:2610.27175v2, 2026.
\bibitem{Gl26} A.~Glazyrin, \href{https://arxiv.org/abs/2608.16116}{New constructions of optimal arrangements of $2d$ lines in $\mathbb C^d$}, arXiv:2608.16116, 2026.
\bibitem{GS67} J.-M.~Goethals and J.~J.~Seidel, Orthogonal matrices with zero diagonal, Canad.\ J.\ Math.\ \textbf{19} (1967), 1001--1010.
\bibitem{KK98} F.~Kamali and H.~Kharaghani, \href{https://ajc.maths.uq.edu.au/pdf/18/ocr-ajc-v18-p139.pdf}{Dihedral Golay sequences}, Australas.\ J.\ Combin.\ \textbf{18} (1998), 139--145.
\bibitem{Kh92} H.~Kharaghani, \href{https://ajc.maths.uq.edu.au/pdf/6/ocr-ajc-v6-p293.pdf}{Block Golay sequences with applications}, Australas.\ J.\ Combin.\ \textbf{6} (1992), 293--303.
\bibitem{Ku26} K.~Kumar, \href{https://github.com/vatskishan03/Hadamard}{New Hadamard matrices} (software repository), \texttt{complex-hadamard/CONFERENCE.md} and certificate \texttt{C1286.json} at commit \texttt{8685428} (October 2026); table of least known orders \texttt{FRONTIER.md} at commit \texttt{d885022} (6 November 2026).
\bibitem{Mi91} M.~Miyamoto, A construction of Hadamard matrices, J.\ Combin.\ Theory Ser.\ A \textbf{57} (1991), 86--108.
\bibitem{SY20} J.~Seberry and M.~Yamada, Hadamard Matrices: Constructions using Number Theory and Linear Algebra, Wiley, 2020.
\bibitem{SY92} J.~Seberry and M.~Yamada, Hadamard matrices, sequences, and block designs, in: J.~H.~Dinitz and D.~R.~Stinson (eds.), Contemporary Design Theory: A Collection of Surveys, Wiley, New York, 1992, 431--560.
\bibitem{Tu70} R.~J.~Turyn, Complex Hadamard matrices, in: Combinatorial Structures and their Applications, Gordon and Breach, New York, 1970, 435--437.
\bibitem{Tu71} R.~J.~Turyn, On C-matrices of arbitrary powers, Canad.\ J.\ Math.\ \textbf{23} (1971), 531--535.
\bibitem{Ya89} M.~Yamada, Some new series of Hadamard matrices, J.\ Austral.\ Math.\ Soc.\ Ser.\ A \textbf{46} (1989), 371--383.
\bibitem{Za26} V.~Zaitsev, \href{https://www.cs.uleth.ca/~zaitsevv/quaternary.html}{Quaternary Hadamard matrices} (web page with matrix files), accessed 24 November 2026; matrix \texttt{QH566.txt}.
\end{thebibliography}
\end{document}
